\section{A local obstruction to the proposed permanent endpoint construction}

\subsection*{Finding F-027: an isolated binary vertex prevents full quadratic jets}
Let \(k\) be an algebraically closed field, let \(m\geq2\),
and use binary variables \(x,y\) and transverse variables
\(z_1,\ldots,z_R\), where \(R\geq2m-1\).
Let
\[
 A(x,y,z)=A_0(x,y)+B(z)
\]
be an \(m\)-by-\(m\) matrix of homogeneous linear forms.
Write \(\pi_{\leq2}\) for projection of a degree-\(m\)
form onto its monomials of total transverse degree at most
two. Consider the full pencil parameter map
\[
 \mathcal P(A)=\pi_{\leq2}\operatorname{per}_m A .
                                                               \tag{F-027.1}
\]
Every coefficient of every matrix entry may vary.

Suppose there are a nonzero binary point
\(p=(p_x,p_y)\in k^2\) and an index \(i\) such that
the entire \(i\)-th row and the entire \(i\)-th column of
\(A_0(p)\) are zero. Then the differential of (F-027.1)
at \(A\) is not surjective. The same assertion holds
with determinant in place of permanent.

\begin{proof}
Consider the linear map \(T:k^2\oplus k^R\to k^{2m-1}\)
whose components are the entries in the union of the
\(i\)-th row and the \(i\)-th column of \(A\).
Its kernel has dimension at least
\[
                   R+2-(2m-1)\geq2,
\]
and it contains \((p,0)\).
The restriction of \(T\) to the binary plane has rank
at most one, because \(p\ne0\) belongs to its kernel.
If that rank is one, its binary kernel is precisely
the line \(kp\). The full kernel therefore contains
a vector \((q,w)\) with \(w\ne0\).
Along the affine line
\[
                     (x,y,z)=(p+tq,tw)
\]
the entire \(i\)-th row and \(i\)-th column of \(A\)
vanish identically in \(t\).

For any entry variation \(H\), the derivative of the
permanent is
\[
 \left.\frac{d}{d\epsilon}\right|_{\epsilon=0}
        \operatorname{per}(A+\epsilon H)
   =\sum_{a,b}H_{ab}\operatorname{per}A_{\widehat a,\widehat b}.
\]
On the displayed line, every cofactor except the one
with \((a,b)=(i,i)\) is zero: deleting any other row or
column retains a zero row or a zero column. Consequently,
if
\[
 P(t)=\operatorname{per}A_{\widehat i,\widehat i}(p+tq,tw),
\]
every derivative restricts to
\[
                    (\alpha+\beta t)P(t).
                                                               \tag{F-027.2}
\]
Here \(\alpha,\beta\in k\) depend on \(H_{ii}\).
After reduction modulo \(t^3\), the right side spans
at most a two-dimensional subspace of \(k[t]/(t^3)\).

In contrast, evaluation on this line followed by reduction
modulo \(t^3\) is surjective from the full target of
(F-027.1) onto \(k[t]/(t^3)\). Indeed choose binary
and transverse linear forms \(L,\zeta\) with
\(L(p)=1\) and \(\zeta(w)=1\). The three target forms
\[
                  L^m,\qquad L^{m-1}\zeta,\qquad
                           L^{m-2}\zeta^2
                                                               \tag{F-027.3}
\]
evaluate respectively to
\[
             (1+at)^m,\qquad (1+at)^{m-1}t,\qquad
                       (1+at)^{m-2}t^2,\qquad a=L(q).
\]
Their coefficient vectors modulo \(t^3\) form a triangular
matrix with diagonal entries one, so they span
\(k[t]/(t^3)\) in every characteristic.
Terms removed by \(\pi_{\leq2}\) evaluate to multiples
of \(t^3\), so evaluation commutes with that projection
modulo \(t^3\).
Thus a surjective differential in (F-027.1) would induce
a surjective map to the three-dimensional space
\(k[t]/(t^3)\), contradicting (F-027.2).

It remains to handle binary rank zero for \(T\).
In that case the entire incident row and column of
\(A_0(x,y)\) vanish identically on the binary plane.
On that plane, the differential of the full permanent
has the form
\[
               H_{ii}(x,y)\,
                    \operatorname{per}(A_0)_{\widehat i,\widehat i}.
\]
Its space of degree-\(m\) binary forms has dimension
at most two, whereas the binary part of the target
has dimension \(m+1\geq3\). Thus the differential is
again not surjective, independently of the transverse
variables.

For the determinant, the derivative has the usual signed
cofactors. Exactly the same cofactors vanish, and the
remaining diagonal cofactor has positive sign, so the
argument is unchanged.
\end{proof}

\paragraph{An explicit annihilator when the complementary permanent is nonzero.}
Write
\[
                    P(t)=P_0+P_1t+P_2t^2+O(t^3).
\]
If \(P_0\ne0\), the nonzero functional
\[
 (c_0,c_1,c_2)\longmapsto
       (P_1^2-P_0P_2)c_0-P_0P_1c_1+P_0^2c_2
                                                               \tag{F-027.4}
\]
annihilates every vector in (F-027.2) modulo \(t^3\).
This makes the missing target direction explicit.
The dimension proof does not require \(P_0\ne0\).

\paragraph{Consequence for the proposed block construction.}
Suppose the binary part is block diagonal with any isolated
one-by-one block \(\ell(x,y)\). Over \(k\), a nonzero
homogeneous binary linear form has a nonzero root \(p\);
if \(\ell=0\), choose any nonzero \(p\).
At that root the associated row and column of \(A_0(p)\)
are zero. F-027 therefore applies independently of every
choice of transverse labels.
In particular, an arbitrary non-diagonal \(3\)-by-\(3\)
binary core followed by a nonempty diagonal remainder
cannot produce a full-rank transverse-order-two Jacobian
in either \(r=2m+1\) or \(r=2m+2\) variables.
The moving-binary-point refinement is essential for
the first of these dimensions.

This is a limitation of the specified base pencils.
It does not disprove dominance of the unrestricted
permanent coefficient map; the finite generic endpoint
witnesses in A-031 are consistent with the result.
The proof leaves open constructions whose binary part
has no isolated vertex, including connected binary
pencils and decompositions into blocks of size at least
two. It gives no such obstruction at \(R=2m-2\),
where the kernel of \(T\) can consist only of the
binary line \(kp\). This is consistent with O-027's
full construction in \(r=2m\) variables.

\paragraph{Geometric dictionary and scope.}
The argument concerns the tangent map of an affine
coefficient morphism and a linear quotient of its target.
Equivalently, substitute \(A+\epsilon H\) over
\(k[\epsilon]/(\epsilon^2)\), and then evaluate the
resulting tangent polynomial in \(k[t]/(t^3)\).
This uses the same explicit finite coefficient-family
framework as M-013: FGA, Bourbaki~221, following
Proposition~3.8, pp.~265--266. The monomial gradings and
coordinate-comodule maps have the precise dictionary
SGA~3, I, Propositions~4.7.2--4.7.3, pp.~34--36.
The dual-number and \(t\)-adic substitutions above are
written out as polynomial identities; no separate
smoothness or lifting theorem is invoked.
Every identity is an integral polynomial identity;
the proof works in every characteristic once the displayed
binary point and transverse kernel vector exist.
It makes no inference about a coordinate ideal or an
irreducible multiplicity from a rank defect at one base
point, and uses no virtual positivity argument.
